% sections/05_commercial.tex -- Q7/Q8/Q9/Q10/Q11 商业分析与近未来展望

\section{商业分析与近未来展望}

\subsection{Q7：轨道vs地面数据中心TCO全生命周期对比}

\subsubsection{地面数据中心TCO基线}

1 MW IT负载的地面AI数据中心10年TCO结构\cite{cloudzero2026gpu,cgranier2025capex,spheron2026gpu}：

\begin{table}[H]
\centering
\caption{地面AI数据中心10年TCO（1 MW IT负载）}
\label{tab:ground_tco}
\small
\begin{tabular}{lrc}
\toprule
\textbf{成本项} & \textbf{金额（US\$M）} & \textbf{占比} \\
\midrule
基建外壳（土地+建筑+电力+冷却） & \$7--12M & 15--20\% \\
GPU硬件（GB300 NVL72，$\sim$5柜/MW） & \$22--28M & 55--65\% \\
网络+存储+其他IT & \$5--8M & 15--20\% \\
\rowcolor{gray!10} \textbf{总Capex} & \textbf{\$35--48M} & \textbf{100\%} \\
电力（10年，\$0.05--0.10/kWh） & \$5--10M & -- \\
冷却能耗（PUE 1.1--1.3） & \$3--5M & -- \\
运维 & \$3--5M & -- \\
退役 & \$1--2M & -- \\
\rowcolor{gray!10} \textbf{10年TCO合计} & \textbf{\$37.8--91.4M} & \textbf{--} \\
\bottomrule
\end{tabular}
\end{table}

\subsubsection{轨道数据中心TCO——多情景对比}

\begin{table}[H]
\centering
\caption{轨道vs地面数据中心10年TCO多情景对比}
\label{tab:orbital_tco}
\begin{tabularx}{\textwidth}{lcccc}
\toprule
\textbf{情景} & \textbf{轨道TCO} & \textbf{地面低端} & \textbf{地面高端} & \textbf{比值} \\
\midrule
乐观（\$50/kg, 溢价2.5x, 寿命7年） & \$261.4M & \$34.6M & \$88.3M & 4.3x \\
基准（\$200/kg, 溢价3.5x, 寿命5年） & \$681.0M & \$37.8M & \$91.4M & 10.5x \\
悲观（\$500/kg, 溢价5x, 寿命3年） & \$1,220M & \$42.1M & \$95.8M & 17.7x \\
\bottomrule
\end{tabularx}
\end{table}

轨道数据中心10年TCO的构成（基准情景）\cite{note_tco_script}：

\begin{table}[H]
\centering
\caption{轨道数据中心10年TCO构成（基准情景 \$200/kg）}
\label{tab:orbital_detail}
\begin{tabular}{lrl}
\toprule
\textbf{成本项} & \textbf{金额（US\$M）} & \textbf{说明} \\
\midrule
发射成本 & \$2.0M & 1次Starship，200t \\
计算硬件（抗辐射加固） & \$34.0M & 溢价3--5x地面 \\
散热系统 & \$3.5M & 辐射散热板+液冷回路 \\
光伏+储能 & \$300.0M & 5,000 m\textsuperscript{2}空间级GaAs \\
卫星平台 & \$30.0M & 姿态控制+结构+通信 \\
硬件替换（设备） & \$257.3M & 寿命5年后替换 \\
硬件替换（发射） & \$1.0M & 替换卫星的发射成本 \\
地面运维 & \$50.0M & 地面站+测控+运维人员 \\
退役 & \$3.0M & 受控再入/坟墓轨道 \\
\rowcolor{gray!10} \textbf{10年TCO合计} & \textbf{\$681.0M} & \textbf{vs地面中位\$64.6M} \\
\bottomrule
\end{tabular}
\end{table}

\begin{figure}[H]
\centering
% figures/tco_breakdown.tex
% TCO 成本构成堆叠柱状图 —— 轨道 vs 地面
\begin{tikzpicture}
\begin{axis}[
    width=0.95\textwidth,
    height=8cm,
    ybar stacked,
    bar width=55pt,
    ylabel={\textbf{10年TCO (US\$\,M)}},
    ymin=0, ymax=750,
    xtick={1,2},
    xticklabels={\textbf{轨道数据中心}\\\textbf{\$681M}, \textbf{地面数据中心}\\\textbf{\$65M}},
    xticklabel style={align=center, font=\footnotesize\bfseries},
    ytick={0,100,200,300,400,500,600,700},
    yticklabel style={font=\footnotesize},
    legend style={
        at={(0.5,-0.25)},
        anchor=north,
        font=\fontsize{6}{8}\selectfont,
        legend columns=4,
        cells={anchor=center}
    },
    enlarge x limits=0.35,
]

% 每个 \addplot 同时给出两柱的值
\addplot[fill=primary!30,  draw=primary!80]  coordinates {(1,300)(2,0)};
\addlegendentry{光伏+储能 \$300M};

\addplot[fill=primary!50,  draw=primary!80]  coordinates {(1,258)(2,0)};
\addlegendentry{硬件替换 \$258M};

\addplot[fill=primary!70,  draw=primary!80]  coordinates {(1,34)(2,39)};
\addlegendentry{计算硬件 \$34M};

\addplot[fill=accent!40,   draw=accent!80]   coordinates {(1,30)(2,10)};
\addlegendentry{卫星平台 \$30M};

\addplot[fill=accent!60,   draw=accent!80]   coordinates {(1,50)(2,26)};
\addlegendentry{地面运维 \$50M};

\addplot[fill=accent!80,   draw=accent!80]   coordinates {(1,2)(2,8)};
\addlegendentry{发射 \$2M};

\addplot[fill=graylight!60,draw=graybase!60] coordinates {(1,7)(2,5)};
\addlegendentry{其他 \$7M};

\addplot[fill=graybase!30, draw=graybase!60] coordinates {(1,0)(2,2)};
\addlegendentry{退役};

% 总量标注
\node[above, font=\footnotesize\bfseries, primary]
    at (axis cs:1,690) {\$681M};
\node[above, font=\footnotesize\bfseries, primary]
    at (axis cs:2,90) {\$65M};

% 右侧柱各段标注
\node[anchor=west, font=\tiny, text=primary!70]
    at (axis cs:2.2,18) {GPU硬件及刷新 \$39M};
\node[anchor=west, font=\tiny, text=accent!40]
    at (axis cs:2.2,33) {基建 \$10M};
\node[anchor=west, font=\tiny, text=accent!60]
    at (axis cs:2.2,53) {网络存储 \$26M};
\node[anchor=west, font=\tiny, text=accent!80]
    at (axis cs:2.2,68) {电力 \$8M};
\node[anchor=west, font=\tiny, text=graybase!60]
    at (axis cs:2.2,75) {运维 \$5M};
\node[anchor=west, font=\tiny, text=graybase!40]
    at (axis cs:2.2,82) {退役 \$2M};

\end{axis}
\end{tikzpicture}

\caption{轨道数据中心 vs 地面数据中心 10年TCO成本构成堆叠对比。左柱：轨道数据中心总TCO \$681M，光伏+储能（\$300M）和硬件替换（\$258M）为最大成本项；右柱：地面数据中心总TCO \$65M，GPU硬件及刷新（\$39M）和网络存储（\$26M）占主导。轨道TCO是地面的$\sim$10.5倍。}
\label{fig:tco_breakdown}
\end{figure}

\subsubsection{马斯克"5年内更便宜"的敏感性分析}

使太空TCO与地面持平所需的临界值\cite{pitchbook2026spacex}：

\begin{table}[H]
\centering
\caption{关键临界值：太空算力TCO与地面持平的条件}
\label{tab:threshold}
\small
\begin{tabular}{lp{2.5cm}p{1.8cm}p{1.5cm}p{3.5cm}}
\toprule
\textbf{参数} & \textbf{当前值} & \textbf{临界值} & \textbf{差距} & \textbf{达成时间估计} \\
\midrule
发射成本（\$/kg） & $\sim$\$500--2,700 & $<$\$100/kg & 5--27x & 乐观2028，基准2030+ \\
轨道硬件成本（\$/W） & $\sim$\$80--200/W & $<$\$20/W & 4--10x & 需大批量生产+标准化 \\
在轨硬件寿命（年） & 估计3--5 & $>$7年 & $\sim$2x & 未知，待Rubin在轨验证 \\
GPU能效（TFLOPS/W） & $\sim$7.5（B200 FP4） & $>$20 & -- & Rubin R100已$\sim$21.7 \\
空间光伏成本（\$/W） & \$200--400 & $<$\$50/W & 4--8x & 需规模化生产 \\
\bottomrule
\end{tabular}
\end{table}

\textbf{条件组合判断}：发射成本$<$\$100/kg + 硬件寿命$>$7年 + 硬件成本溢价$<$3x地面 $\rightarrow$ 轨道TCO为\$408.1M（6.3x地面），即使在翻转条件组合下仍未达到地面水平。乐观估计2029--2032年可达临界可行（概率$<$25\%）。

\subsection{Q8：GPU/ASIC能效的历史演进与近未来外推}

\subsubsection{NVIDIA GPU架构TFLOPS/W历史曲线}

\begin{table}[H]
\centering
\caption{NVIDIA数据中心旗舰GPU能效演进（2012--2026）}
\label{tab:gpu_efficiency}
\footnotesize
\begin{tabularx}{\textwidth}{lccccccr}
\toprule
\textbf{架构} & \textbf{年份} & \textbf{旗舰} & \textbf{制程} & \textbf{TDP (W)} & \textbf{AI TFLOPS} & \textbf{AI TFLOPS/W} & \textbf{精度} \\
\midrule
Kepler & 2012 & K40 & 28nm & 235 & -- & -- & -- \\
Pascal & 2016 & P100 & 16nm & 300 & 21.2 & 0.071 & FP16 \\
Volta & 2017 & V100 & 12nm & 300 & 125 & 0.42 & FP16 Tensor \\
Ampere & 2020 & A100 & 7nm & 400 & 312 & 0.78 & FP16 Tensor \\
Hopper & 2022 & H100 & 4nm & 700 & 1,979 & 2.83 & FP8 Tensor \\
Blackwell & 2024 & B200 & 4NP & 1,200 & $\sim$9,000 & 7.5 & FP4 Tensor \\
\textbf{Rubin} & \textbf{2026} & \textbf{R100} & \textbf{3nm} & \textbf{$\sim$2,300} & \textbf{$\sim$50,000} & \textbf{$\sim$21.7} & \textbf{FP4 Tensor} \\
\bottomrule
\end{tabularx}
\end{table}

\textbf{关键趋势}：AI精度能效（FP16$\rightarrow$FP4 Tensor）8年提升约52倍，每代约2.5--3x提升——从Volta的0.42攀升至Rubin的$\sim$21.7。但绝对功耗同步飙升10倍（235~W$\rightarrow$2{,}300~W）
\cite{nvidia2026rubin,spheron2026rubin,barrack2026rubin}。

\begin{figure}[H]
\centering
% figures/gpu_paradox.tex
% GPU 能效悖论：AI TFLOPS/W vs TDP 双 y 轴
\begin{tikzpicture}
\begin{axis}[
    width=0.95\textwidth,
    height=7cm,
    xlabel={\textbf{年份 / 架构}},
    ylabel={\textbf{AI TFLOPS/W}（对数刻度）},
    ymode=log,
    log basis y={10},
    ymin=0.03, ymax=30,
    ytick={0.05,0.1,0.2,0.5,1,2,5,10,20},
    yticklabels={0.05,0.1,0.2,0.5,1,2,5,10,20},
    ylabel style={color=primary},
    tick label style={font=\footnotesize},
    xtick={1,2,3,4,5,6,7},
    xticklabels={2012\\Kepler, 2016\\Pascal, 2017\\Volta, 2020\\Ampere, 2022\\Hopper, 2024\\Blackwell, 2026\\Rubin},
    xticklabel style={align=center, font=\footnotesize},
    xmin=0.5, xmax=7.5,
    grid=major,
    legend style={
        at={(0.02,0.98)},
        anchor=north west,
        font=\footnotesize,
        cells={anchor=west}
    },
]

% 左 y 轴：AI TFLOPS/W（蓝色）
\addplot[color=primary, line width=1.5pt, mark=*, mark size=3pt]
    coordinates {(2,0.071)(3,0.42)(4,0.78)(5,2.83)(6,7.5)(7,21.7)};
\addlegendentry{AI TFLOPS/W (FP Tensor)}

% 右 y 轴：TDP (W)（红色）
\end{axis}

\begin{axis}[
    width=0.95\textwidth,
    height=7cm,
    axis x line=none,
    axis y line*=right,
    ylabel={\textbf{TDP (W)}},
    ylabel style={color=highlight},
    ymin=0, ymax=2600,
    ytick={0,500,1000,1500,2000,2500},
    yticklabel style={font=\footnotesize, color=highlight},
    xmin=0.5, xmax=7.5,
    legend style={
        at={(0.02,0.88)},
        anchor=north west,
        font=\footnotesize,
        cells={anchor=west}
    },
]

\addplot[color=highlight, line width=1.5pt, dashed, mark=square*, mark size=3pt]
    coordinates {(1,235)(2,300)(3,300)(4,400)(5,700)(6,1200)(7,2300)};
\addlegendentry{TDP (W)}

% 标注：能效 vs 功耗悖论
\draw[->, >=stealth, graybase, line width=0.7pt]
    (axis cs:4,1800) -- (axis cs:6.5,2000);
\node[graybase, font=\footnotesize\bfseries, anchor=south, align=center]
    at (axis cs:5,1950) {\textbf{能效 $\uparrow$52x}\\\textbf{但功耗 $\uparrow$10x}};

\end{axis}
\end{tikzpicture}

\caption{GPU能效悖论：AI TFLOPS/W（左轴，对数刻度，蓝色实线）持续提升$\sim$52倍，但TDP（右轴，红色虚线）同步飙升$\sim$10倍。能效改善每瓦产出，却不能降低每芯片的散热需求——Dennard缩放已失效，对轨道数据中心的散热瓶颈缓解有限。}
\label{fig:gpu_paradox}
\end{figure}

\subsubsection{2030年GPU能效推算与热密度影响}

基于台积电工艺路线图（N3$\rightarrow$N2$\rightarrow$A14 1.4nm$\rightarrow$A13/A12 1.2--1.3nm）\cite{tsmc2026symposium,fiisual2026tsmc}，2030年GPU能效预计达$\sim$100--150 TFLOPS/W（FP4），相同算力所需GPU数量减少5--7x。

但\textbf{芯片级热密度仍在增加}（Dennard缩放已失效），轨道散热面积需求与算力几乎成正比。能效提升改善的是每瓦产出而非每芯片散热需求——\textbf{对轨道数据中心的散热瓶颈缓解有限}。

\subsection{Q9：发射成本的历史下降曲线与近未来外推}

\subsubsection{Wright's Law适用性}

SpaceX发射成本15年间下降约100倍（Shuttle \$54,500/kg$\rightarrow$Starship目标$\sim$\$100--200/kg），遵循航空航天领域标准85\%学习率（累计产量翻倍时单位成本降15\%）\cite{nasalaunch2018}。

\textbf{局限性}：Wright's Law仅适用于制造成本。33fg Research（2026）指出，\$/kg渐进线在$\sim$\$10/kg（制造成本极限）+ 燃料/运营$\sim$\$10--20/kg $\rightarrow$ \textbf{物理地板约\$20--30/kg}\cite{33fg2026starship}。

\subsubsection{2030年弹性预估}

\begin{table}[H]
\centering
\caption{发射成本2030年弹性预估}
\label{tab:launch_2030}
\small
\begin{tabular}{lp{1.7cm}p{2.2cm}p{5.5cm}p{2.5cm}}
\toprule
\textbf{情景} & \textbf{2030 \$/kg} & \textbf{年发射频次} & \textbf{关键假设} & \textbf{概率评估} \\
\midrule
乐观 & \$50/kg & 200+/年 & FAA大幅放宽、复用$>$30次/艘、年产$>$500艘 & 低（$<$20\%） \\
基准 & \$200/kg & 100--150/年 & FAA扩展至100+/年、Block 3成熟、复用10--20次 & 中（$\sim$50\%） \\
悲观 & \$500/kg & 50--70/年 & FAA维持69/年上限、复用率$<$5次 & 中偏低（$\sim$30\%） \\
\bottomrule
\end{tabular}
\end{table}

\textbf{核心洞察}：即使发射成本降至\$50/kg（乐观），仅此一项不足以使太空TCO低于地面——硬件成本溢价和寿命短板是同等重要的变量\cite{spaceodyssey2026,exterra2026}。

\subsection{Q10：下一代散热材料与结构设计对辐射散热极限的突破潜力}

\subsubsection{关键技术路线图}

\begin{table}[H]
\centering
\caption{下一代空间散热技术路线图}
\label{tab:thermal_roadmap}
\small
\begin{tabular}{p{3.5cm}clp{4.0cm}}
\toprule
\textbf{技术方向} & \textbf{TRL} & \textbf{预计部署} & \textbf{关键参数改进} \\
\midrule
石墨烯涂层热管 & 4--5 & 2028--2032 & 热导率$>$5,000 W/(m·K)\cite{esa2017graphene,scalia2023graphene} \\
\hline
增材制造钛环路热管+可展开辐射板 & 5--6 & 2027--2030 & 热传输$>$10m，面积$\times$2--3\cite{nasa2024tfaws} \\
\hline
形状记忆合金自展开辐射器 & 3--4 & 2029--2033 & 被动热驱动展开\cite{pennstate2024sma} \\
\hline
FENNEC折叠式CubeSat辐射器 & 5--6 & 2026--2028 & CubeSat散热能力$\times$2--3\cite{esa2024fennec} \\
\hline
微结构表面增强辐射率 & 3--5 & 2030--2035 & $\varepsilon$：0.85$\rightarrow$0.95+ \\
\hline
空间两相流泵驱动冷却 & 5--6 & 2026--2030 & 等效热导率百倍于固体 \\
\hline
折纸式/充气式散热面 & 2--4 & 2032+ & 超高收纳比（$>$10:1） \\
\bottomrule
\end{tabular}
\end{table}

\subsubsection{散热密度天花板突破推演}

基于Q3（真实热环境惩罚）和Q10（材料/结构进步）的综合修正\cite{nasa2024sst}：

\begin{table}[H]
\centering
\caption{散热密度进化路径与1 GW散热面积修正}
\label{tab:flux_evolution}
\small
\begin{tabular}{p{3.0cm}cccc}
\toprule
\textbf{技术组合} & $\varepsilon$ & $T$ (\unitC{}) & \textbf{通量 (W/m\textsuperscript{2})} & \textbf{1 GW面积 (km\textsuperscript{2})} \\
\midrule
\textbf{AI1当前（BOL）} & 0.90 & 69.5 & 1,400 & 0.71 \\
Level 1：高$\varepsilon$涂层+80\unitC{} & 0.93 & 80 & 1,740 & 0.57 \\
Level 2：+微结构$\varepsilon$=0.96+90\unitC{} & 0.96 & 90 & 2,150 & 0.47 \\
Level 3：+Ti-石墨辐射翼+120\unitC{} & 0.96 & 120 & 2,970 & 0.34 \\
Level 4：+折纸展开$\times$5面积 & 0.96 & 100 & $\sim$2,130$\rightarrow\geq$等效10,650 & $\sim$0.09 \\
\rowcolor{gray!10}
\textbf{综合2030--32最现实估计} & \textbf{0.93} & \textbf{75--100} & \textbf{1,500--2,000} & \textbf{0.50--0.67} \\
\bottomrule
\end{tabular}
\end{table}

\textbf{结论}：散热密度翻倍（1,400$\rightarrow$2,800+ W/m\textsuperscript{2}）在2035年前的单一技术路径上不可行，但"高温芯片 + 耐久涂层 + 适度展开"的组合路径可在2032年前实现40--50\%改善\cite{nasatfaws2024ti,act2026deployable,kanti2025graphene}。

\subsection{Q11："静态今天"vs"近未来外推"对比元分析}

\subsubsection{结论对比表}

\begin{table}[H]
\centering
\caption{静态（2026）vs近未来（2030）结论对比}
\label{tab:static_vs_future}
\small
\begin{tabular}{p{2.5cm}p{4.5cm}p{4.5cm}}
\toprule
\textbf{维度} & \textbf{静态结论（2026）} & \textbf{近未来结论（2030）} \\
\midrule
发射成本 & \$500--2,700/kg & \$50--200/kg（Starship复用+学习曲线） \\
GPU能效 & B200 $\sim$7.5 TFLOPS/W & $\sim$100--150 TFLOPS/W（TSMC 3nm$\rightarrow$1.4nm） \\
散热物理 & $\sim$1,400\unitWpm{}天花板 & 1,500--2,000\unitWpm{}（材料进步20--50\%） \\
太阳能供电 & 30\%效率GaAs，1MW需$\sim$5,000 m\textsuperscript{2} & 36\%效率五结电池，1MW需$\sim$4,000 m\textsuperscript{2} \\
硬件可靠性 & 轨道寿命3--5年（估计） & 可能5--7年（Rubin在轨验证+设计迭代） \\
地面竞争 & 地面TCO \$37.8--91.4M/10年/MW（中位\$64.6M） & 地面也在进步（不是静止参照系） \\
\rowcolor{gray!10}
\textbf{总体判断} & \textbf{TCO是地面的10.5倍，不可行} & \textbf{乐观下可接近4.3x，仍显著高于地面} \\
\bottomrule
\end{tabular}
\end{table}

\subsubsection{关键阈值矩阵：翻转条件组合}

\begin{table}[H]
\centering
\caption{太空算力从"不可行"翻转为"可行"的条件组合}
\label{tab:flip_matrix}
\footnotesize
\begin{tabular}{p{1.3cm}p{1.6cm}p{1.5cm}p{1.3cm}p{1.2cm}p{1.1cm}p{1.3cm}p{1.4cm}}
\toprule
\textbf{组合} & \textbf{发射成本} & \textbf{GPU能效} & \textbf{散热密度} & \textbf{硬件寿命} & \textbf{硬件溢价} & \textbf{预计时间} & \textbf{概率} \\
\midrule
A：最乐观 & $<$\$50/kg & $>$150 TF/W & $>$2,500 & $>$7年 & $<$2x & 2032+ & 极低（$<$10\%） \\
B：乐观可行 & $<$\$100/kg & $>$100 TF/W & $>$2,000 & $>$5年 & $<$3x & 2030--32 & 低（15--25\%） \\
C：临界可行 & $<$\$200/kg & $>$80 TF/W & $>$1,800 & $>$5年 & $<$4x & 2032--35 & 中低（25--40\%） \\
D：仍不可行 & $>$\$200/kg & $<$80 TF/W & $<$1,800 & $<$5年 & $>$4x & 2035+ & 当前状态 \\
\bottomrule
\end{tabular}
\end{table}

\begin{figure}[H]
\centering
% figures/threshold_radar.tex
% 临界阈值蛛网图：四轴雷达图（发射成本/GPU能效/散热密度/硬件寿命）
% 使用 TikZ 手工绘制四边形雷达图

\begin{tikzpicture}[scale=1.0]

% ===== 轴定义 =====
% 四个轴：右(0°)、上(90°)、左(180°)、下(270°)
% 中心点 (0,0)，最大半径 3.5cm

\def\rmax{3.5}

% ===== 网格线（0.25, 0.5, 0.75, 1.0） =====
\foreach \r in {0.875,1.75,2.625,3.5}{
    \draw[gray!30, line width=0.3pt]
        (\r,0) -- (0,\r) -- (-\r,0) -- (0,-\r) -- cycle;
}

% ===== 四轴 =====
% 轴 1: 右 — GPU 能效 (TFLOPS/W)
% 轴 2: 上 — 散热密度 (W/m²)
% 轴 3: 左 — 发射成本 ($/kg, 反向)
% 轴 4: 下 — 硬件寿命 (年)
\draw[->, >=stealth, graybase, line width=0.8pt] (0,0) -- ( \rmax,0);
\draw[->, >=stealth, graybase, line width=0.8pt] (0,0) -- (0, \rmax);
\draw[->, >=stealth, graybase, line width=0.8pt] (0,0) -- (-\rmax,0);
\draw[->, >=stealth, graybase, line width=0.8pt] (0,0) -- (0,-\rmax);

% ===== 轴标签 =====
\node[primary, font=\bfseries\footnotesize, anchor=west]
    at (\rmax+0.3,0) {\textbf{GPU能效}};
\node[primary, font=\scriptsize, anchor=west]
    at (\rmax+0.3,-0.35) {TFLOPS/W};

\node[primary, font=\bfseries\footnotesize, anchor=south]
    at (0,\rmax+0.3) {\textbf{散热密度}};
\node[primary, font=\scriptsize, anchor=south]
    at (0,\rmax+0.05) {W/m\textsuperscript{2}};

\node[primary, font=\bfseries\footnotesize, anchor=east]
    at (-\rmax-0.3,0) {\textbf{发射成本}};
\node[primary, font=\scriptsize, anchor=east]
    at (-\rmax-0.3,-0.35) {\$/kg (反向)};

\node[primary, font=\bfseries\footnotesize, anchor=north]
    at (0,-\rmax-0.3) {\textbf{硬件寿命}};
\node[primary, font=\scriptsize, anchor=north]
    at (0,-\rmax-0.05) {年};

% ===== 刻度标记 =====
% GPU能效: 0, 25, 50, 75, 100, 125, 150
\foreach \v/\l in {0.583/25, 1.167/50, 1.75/75, 2.333/100, 2.917/125, 3.5/150}{
    \draw[gray!40] (\v,0.08) -- (\v,-0.08);
    \node[font=\tiny, graybase, anchor=north] at (\v,-0.1) {\l};
}

% 散热密度: 0, 500, 1000, 1500, 2000, 2500, 3000
\foreach \v/\l in {0.583/500, 1.167/1000, 1.75/1500, 2.333/2000, 2.917/2500, 3.5/3000}{
    \draw[gray!40] (0.08,\v) -- (-0.08,\v);
    \node[font=\tiny, graybase, anchor=east] at (-0.12,\v) {\l};
}

% 发射成本 (反向): 0, 50, 100, 150, 200, 250, 300
\foreach \v/\l in {0.583/50, 1.167/100, 1.75/150, 2.333/200, 2.917/250, 3.5/300}{
    \draw[gray!40] (-\v,0.08) -- (-\v,-0.08);
    \node[font=\tiny, graybase, anchor=north] at (-\v,-0.1) {\l};
}

% 硬件寿命: 0, 2, 4, 6, 8, 10, 12
\foreach \v/\l in {0.583/2, 1.167/4, 1.75/6, 2.333/8, 2.917/10, 3.5/12}{
    \draw[gray!40] (0.08,-\v) -- (-0.08,-\v);
    \node[font=\tiny, graybase, anchor=east] at (-0.12,-\v) {\l};
}

% ===== 红色虚线多边形: 2026 基准 =====
% 发射 $200 → 左轴: 200/300 * 3.5 = 2.333
% GPU能效 7.5 → 右轴: 7.5/150 * 3.5 = 0.175
% 散热 1400 → 上轴: 1400/3000 * 3.5 = 1.633
% 寿命 5 → 下轴: 5/12 * 3.5 = 1.458
\fill[highlight, opacity=0.15]
    (0.175,0) -- (0,1.633) -- (-2.333,0) -- (0,-1.458) -- cycle;
\draw[highlight, line width=1.2pt, dashed]
    (0.175,0) -- (0,1.633) -- (-2.333,0) -- (0,-1.458) -- cycle;
\node[highlight, font=\footnotesize\bfseries, anchor=south west]
    at (0.2,0.2) {\textbf{2026 基准}};

% 数据点标注
\fill[highlight] (0.175,0) circle (1.8pt);
\fill[highlight] (0,1.633) circle (1.8pt);
\fill[highlight] (-2.333,0) circle (1.8pt);
\fill[highlight] (0,-1.458) circle (1.8pt);

% ===== 蓝色实线多边形: 翻转条件 =====
% 发射 <$100 → 100/300 * 3.5 = 1.167
% GPU能效 >100 → 100/150 * 3.5 = 2.333
% 散热 >2000 → 2000/3000 * 3.5 = 2.333
% 寿命 >7 → 7/12 * 3.5 = 2.042
\fill[primary, opacity=0.25]
    (2.333,0) -- (0,2.333) -- (-1.167,0) -- (0,-2.042) -- cycle;
\draw[primary, line width=1.5pt]
    (2.333,0) -- (0,2.333) -- (-1.167,0) -- (0,-2.042) -- cycle;
\node[primary, font=\footnotesize\bfseries, anchor=south east]
    at (1.5,1.5) {\textbf{翻转条件}};

% 数据点标注
\fill[primary] (2.333,0) circle (2pt);
\fill[primary] (0,2.333) circle (2pt);
\fill[primary] (-1.167,0) circle (2pt);
\fill[primary] (0,-2.042) circle (2pt);

% ===== 面积差标注 =====
\node[graybase, font=\tiny\itshape, align=center]
    at (0.8,0.6) {\textbf{面积差}\\\textbf{$=$ 不可行区间}};

% ===== 图例 =====
\node[font=\footnotesize, anchor=north west]
    at (-3.5,3.5) {
    \begin{tabular}{ll}
    \tikz{\draw[highlight, dashed, line width=1.2pt] (0,0) -- (0.5,0);} & 2026 基准 \\
    \tikz{\draw[primary, line width=1.5pt] (0,0) -- (0.5,0);} & 翻转条件 \\
    \end{tabular}
};

\end{tikzpicture}

\caption{临界阈值蛛网图：四轴雷达图对比2026基准（红色虚线）与翻转条件（蓝色实线）。四轴分别为GPU能效（TFLOPS/W）、散热密度（W/m\textsuperscript{2}）、发射成本（\$/kg，反向）和硬件寿命（年）。两多边形面积差代表当前不可行区间——需要四个参数同时达到临界值才可能翻转，单一参数突破不足以改变结论。}
\label{fig:threshold_radar}
\end{figure}

\subsubsection{黄仁勋vs马斯克判断的独立检验}

\begin{table}[H]
\centering
\caption{黄仁勋"务实保留"vs马斯克"5年更便宜"的交叉验证}
\label{tab:huang_vs_musk}
\small
\begin{tabular}{p{2.0cm}p{4.0cm}p{3.5cm}p{4.0cm}}
\toprule
\textbf{判断} & \textbf{黄仁勋立场} & \textbf{马斯克立场} & \textbf{本报告独立判断} \\
\midrule
散热是瓶颈？ & 是（"需要极大散热面积"） & 部分承认但乐观 & \textbf{同意黄仁勋：不受摩尔定律驱动} \\
\hline
5年内成本交叉？ & 不给具体年份 & 明确"4--5年" & \textbf{乐观下可能但概率$<$25\%} \\
\hline
硬件可靠性？ & 隐忧 & 未充分讨论 & \textbf{数据严重不足，关键不确定性} \\
\hline
当前经济性？ & "不理想" & 未直接否认 & \textbf{同意黄仁勋：10.5x地面} \\
\bottomrule
\end{tabular}
\end{table}

\textbf{核心洞察}：没有任何单一参数的突破能使太空算力翻转——需要发射成本、GPU能效、散热密度、硬件寿命\textbf{四个参数同时}达到临界值。这解释了为什么黄仁勋（从芯片工程视角）比马斯克（从文明能源视角）更保守：黄仁勋看到了多参数耦合的工程难度，马斯克假设所有参数会同步改善。

% end of 05_commercial
